%%%%%%%%%%%%%%%%%%%%%%%%%%%%%%%%%%%%%%%%%%%%%%%%%%%%%%%%%%%%%%%%%%%%%%%%%%%%%%
%
% 05_sheaf.tex
%
%%%%%%%%%%%%%%%%%%%%%%%%%%%%%%%%%%%%%%%%%%%%%%%%%%%%%%%%%%%%%%%%%%%%%%%%%%%%%%

\begin{mosbox}{4. Cellular Sheaves: Attaching Local Information}

\BodyFont

The simplicial complex specifies only the \emph{combinatorial structure}
of memory. To represent mathematical information associated with each
concept and relationship, MOS equips the complex with a
\textbf{cellular sheaf}.

A cellular sheaf assigns a vector space to every simplex together with
compatible restriction maps describing how information propagates between
adjacent simplices.

Formally,

\[
\boxed{
\mathcal{F} :
C
\longrightarrow
\mathbf{Vect}
}
\]

where every simplex $\sigma\in C$ receives a stalk

\[
\mathcal F(\sigma)
\]

and every face inclusion

\[
\sigma\subseteq\tau
\]

induces a linear restriction map

\[
\rho_{\tau\rightarrow\sigma}:
\mathcal F(\tau)
\longrightarrow
\mathcal F(\sigma).
\]

The sheaf therefore provides the local algebraic structure that transforms
the simplicial complex into an information-bearing mathematical object.

\end{mosbox}

%%%%%%%%%%%%%%%%%%%%%%%%%%%%%%%%%%%%%%%%%%%%%%%%%%%%%%%%%%%%%%%%%%%%%%%%%%%%%%
%% SHEAF DIAGRAM
%%%%%%%%%%%%%%%%%%%%%%%%%%%%%%%%%%%%%%%%%%%%%%%%%%%%%%%%%%%%%%%%%%%%%%%%%%%%%%

\vspace{5mm}

\begin{center}
\begin{tikzpicture}[
  scale=1.2,
  node/.style={circle, fill=MOSRed, draw=MOSGold, line width=1pt, minimum size=10pt, inner sep=0pt},
  edge/.style={draw=MOSRed, line width=2pt},
  stalk/.style={draw=MOSGold, line width=1.5pt, fill=black, minimum width=2.5cm, minimum height=1cm, text=white, font=\Large\bfseries, rounded corners}
]

% Base space (1-simplex)
\node[node] (v1) at (-3, 0) {};
\node[node] (v2) at (3, 0) {};
\draw[edge] (v1) -- (v2) node[midway, below=3mm, white, font=\Large] {Edge $e$};

\node[white, font=\Large, below=3mm of v1] {Vertex $u$};
\node[white, font=\Large, below=3mm of v2] {Vertex $v$};

% Stalks (Vector Spaces)
\node[stalk] (Fv1) at (-3, 3) {$\mathcal{F}(u)$};
\node[stalk] (Fv2) at (3, 3) {$\mathcal{F}(v)$};
\node[stalk] (Fe)  at (0, 5) {$\mathcal{F}(e)$};

% Restriction Maps
\draw[-{Latex[length=4mm]}, MOSGold, line width=2pt] (Fe) -- (Fv1) node[midway, above left=1mm] {$\rho_{e \to u}$};
\draw[-{Latex[length=4mm]}, MOSGold, line width=2pt] (Fe) -- (Fv2) node[midway, above right=1mm] {$\rho_{e \to v}$};

% Projection lines showing fiber over base
\draw[dashed, MOSRed, line width=1pt] (Fv1) -- (v1);
\draw[dashed, MOSRed, line width=1pt] (Fv2) -- (v2);
\draw[dashed, MOSRed, line width=1pt] (Fe) -- (0,0);

\end{tikzpicture}
\end{center}

%%%%%%%%%%%%%%%%%%%%%%%%%%%%%%%%%%%%%%%%%%%%%%%%%%%%%%%%%%%%%%%%%%%%%%%%%%%%%%
%% LOCAL COMPATIBILITY
%%%%%%%%%%%%%%%%%%%%%%%%%%%%%%%%%%%%%%%%%%%%%%%%%%%%%%%%%%%%%%%%%%%%%%%%%%%%%%

\begin{highlightbox}

\SectionTitle{Local Compatibility}

\BodyFont

The restriction maps encode how information attached to higher-dimensional
relationships must agree with information attached to their faces.

Consequently,

\[
\boxed{
\text{Every local assignment is constrained by its neighbors.}
}
\]

Compatibility is therefore not imposed globally; it is enforced through
the local structure of the sheaf itself.

\end{highlightbox}

%%%%%%%%%%%%%%%%%%%%%%%%%%%%%%%%%%%%%%%%%%%%%%%%%%%%%%%%%%%%%%%%%%%%%%%%%%%%%%
%% GLOBAL SECTION
%%%%%%%%%%%%%%%%%%%%%%%%%%%%%%%%%%%%%%%%%%%%%%%%%%%%%%%%%%%%%%%%%%%%%%%%%%%%%%

\begin{eqbox}

\[
\boxed{
s\in H^{0}(C,\mathcal F)
}
\]

\BodyFont

A \textbf{global section} is a compatible assignment of data to every
simplex satisfying all restriction maps simultaneously.

Within MOS, a global section represents a mathematically coherent memory
state whose local pieces agree everywhere on the simplicial complex.

\end{eqbox}

%%%%%%%%%%%%%%%%%%%%%%%%%%%%%%%%%%%%%%%%%%%%%%%%%%%%%%%%%%%%%%%%%%%%%%%%%%%%%%
%% WHY SHEAVES?
%%%%%%%%%%%%%%%%%%%%%%%%%%%%%%%%%%%%%%%%%%%%%%%%%%%%%%%%%%%%%%%%%%%%%%%%%%%%%%

\begin{remarkbox}

\BodyFont

The cellular sheaf is the mathematical mechanism that links local and
global information.

Rather than storing concepts independently, MOS models memory as a network
of compatible local assignments whose consistency can later be analyzed
using sheaf Laplacians, cohomology and topological diagnostics.

This transition—from combinatorial geometry to algebraic structure—is the
key step that enables the coherence measures developed in the following
sections.

\end{remarkbox}